% ============================================
% DAFTAR SIMBOL
% File: Simbol.tex
% ============================================

% Gunakan package glossaries untuk simbol
% Pastikan sudah ada \makenomenclature atau \makeglossaries di preamble

% ============================================
% SIMBOL HIMPUNAN & STRUKTUR DATA
% ============================================

\newglossaryentry{sym:calT}{
    type=symbols,
    name={\ensuremath{\mathcal{T}}},
    description={Teks input atau dokumen yang akan diproses},
    sort=T
}

\newglossaryentry{sym:calS}{
    type=symbols,
    name={\ensuremath{\mathcal{S}}},
    description={Representasi terstruktur hasil ekstraksi informasi},
    sort=S
}

\newglossaryentry{sym:E}{
    type=symbols,
    name={\ensuremath{E}},
    description={Himpunan entitas dalam \textit{Knowledge Graph}},
    sort=E
}

\newglossaryentry{sym:R}{
    type=symbols,
    name={\ensuremath{R}},
    description={Himpunan relasi dalam \textit{Knowledge Graph}},
    sort=R
}

\newglossaryentry{sym:T}{
    type=symbols,
    name={\ensuremath{T}},
    description={Himpunan \textit{triplet} dalam \textit{Knowledge Graph}},
    sort=T
}

\newglossaryentry{sym:G}{
    type=symbols,
    name={\ensuremath{G}},
    description={\textit{Knowledge Graph} yang didefinisikan sebagai $G = (E, R, T)$},
    sort=G
}

\newglossaryentry{sym:calG}{
    type=symbols,
    name={\ensuremath{\mathcal{G}}},
    description={Graf pengetahuan (\textit{Knowledge Graph}) secara keseluruhan},
    sort=G
}

\newglossaryentry{sym:A}{
    type=symbols,
    name={\ensuremath{A}},
    description={Himpunan seluruh kemungkinan respons/jawaban yang dapat dihasilkan},
    sort=A
}

% ============================================
% SIMBOL ENTITAS & TRIPLET
% ============================================

\newglossaryentry{sym:ei}{
    type=symbols,
    name={\ensuremath{e_i}},
    description={Entitas subjek (\textit{head}) dalam sebuah \textit{triplet}},
    sort=ei
}

\newglossaryentry{sym:ej}{
    type=symbols,
    name={\ensuremath{e_j}},
    description={Entitas objek (\textit{tail}) dalam sebuah \textit{triplet}},
    sort=ej
}

\newglossaryentry{sym:rij}{
    type=symbols,
    name={\ensuremath{r_{ij}}},
    description={Relasi atau predikat yang menghubungkan entitas $e_i$ dan $e_j$},
    sort=rij
}

\newglossaryentry{sym:h}{
    type=symbols,
    name={\ensuremath{h}},
    description={Entitas asal (\textit{head}) dalam format triplet $(h, r, t)$},
    sort=h
}

\newglossaryentry{sym:r}{
    type=symbols,
    name={\ensuremath{r}},
    description={Relasi dalam format triplet $(h, r, t)$},
    sort=r
}

\newglossaryentry{sym:t}{
    type=symbols,
    name={\ensuremath{t}},
    description={Entitas tujuan (\textit{tail}) dalam format triplet $(h, r, t)$},
    sort=t
}

\newglossaryentry{sym:e1}{
    type=symbols,
    name={\ensuremath{e_1}},
    description={\textit{Node} pertama dalam contoh graf sederhana},
    sort=e1
}

\newglossaryentry{sym:e2}{
    type=symbols,
    name={\ensuremath{e_2}},
    description={\textit{Node} kedua dalam contoh graf sederhana},
    sort=e2
}

\newglossaryentry{sym:r1}{
    type=symbols,
    name={\ensuremath{r_1}},
    description={Relasi pertama yang menghubungkan $e_1$ ke $e_2$},
    sort=r1
}

% ============================================
% SIMBOL LLM & GENERASI TEKS
% ============================================

\newglossaryentry{sym:Ptheta}{
    type=symbols,
    name={\ensuremath{P_\theta}},
    description={Fungsi probabilitas yang dipelajari oleh model bahasa dengan parameter $\theta$},
    sort=Ptheta
}

\newglossaryentry{sym:x}{
    type=symbols,
    name={\ensuremath{x}},
    description={Input teks dari pengguna ke model bahasa},
    sort=x
}

\newglossaryentry{sym:yt}{
    type=symbols,
    name={\ensuremath{y_t}},
    description={Token keluaran pada posisi ke-$t$ yang dihasilkan model},
    sort=yt
}

\newglossaryentry{sym:y1yt}{
    type=symbols,
    name={\ensuremath{y_1, \cdots, y_t}},
    description={Urutan token keluaran dari posisi pertama hingga ke-$t$},
    sort=y1yt
}

\newglossaryentry{sym:theta}{
    type=symbols,
    name={\ensuremath{\theta}},
    description={Parameter model \textit{retriever} yang dapat dilatih},
    sort=theta
}

\newglossaryentry{sym:phi}{
    type=symbols,
    name={\ensuremath{\phi}},
    description={Parameter model \textit{generator} yang dapat dilatih},
    sort=phi
}

% ============================================
% SIMBOL GRAPH-RAG
% ============================================

\newglossaryentry{sym:astar}{
    type=symbols,
    name={\ensuremath{a^*}},
    description={Jawaban optimal yang dihasilkan sistem \textit{Graph-RAG}},
    sort=astar
}

\newglossaryentry{sym:a}{
    type=symbols,
    name={\ensuremath{a}},
    description={Jawaban atau respons yang dihasilkan sistem},
    sort=a
}

\newglossaryentry{sym:q}{
    type=symbols,
    name={\ensuremath{q}},
    description={Pertanyaan atau kueri dari pengguna},
    sort=q
}

\newglossaryentry{sym:Q}{
    type=symbols,
    name={\ensuremath{Q}},
    description={Pertanyaan pengguna dalam notasi proses ekstraksi kata kunci},
    sort=Q
}

\newglossaryentry{sym:paq}{
    type=symbols,
    name={\ensuremath{p(a \mid q, \mathcal{G})}},
    description={Distribusi probabilitas jawaban $a$ diberikan pertanyaan $q$ dan graf $\mathcal{G}$},
    sort=paq
}

\newglossaryentry{sym:ptheta}{
    type=symbols,
    name={\ensuremath{p_\theta(G \mid q, \mathcal{G})}},
    description={Model \textit{graph retriever} untuk memilih subgraf relevan},
    sort=ptheta
}

\newglossaryentry{sym:pphi}{
    type=symbols,
    name={\ensuremath{p_\phi(a \mid q, G)}},
    description={Model \textit{answer generator} untuk menghasilkan jawaban dari subgraf},
    sort=pphi
}

\newglossaryentry{sym:Gstar}{
    type=symbols,
    name={\ensuremath{G^*}},
    description={Subgraf optimal yang dipilih oleh \textit{retriever}},
    sort=Gstar
}

% ============================================
% SIMBOL KATA KUNCI & RETRIEVAL
% ============================================

\newglossaryentry{sym:Kclues}{
    type=symbols,
    name={\ensuremath{\mathcal{K}^{clues}}},
    description={Kata kunci petunjuk awal dari basis data vektor},
    sort=Kclues
}

\newglossaryentry{sym:Kl}{
    type=symbols,
    name={\ensuremath{\mathcal{K}_l}},
    description={Kata kunci tingkat rendah (\textit{low-level}) untuk pencarian detail spesifik},
    sort=Kl
}

\newglossaryentry{sym:Kh}{
    type=symbols,
    name={\ensuremath{\mathcal{K}_h}},
    description={Kata kunci tingkat tinggi (\textit{high-level}) untuk konteks tema besar},
    sort=Kh
}

\newglossaryentry{sym:LLMkeyword}{
    type=symbols,
    name={\ensuremath{\text{LLM}_{keyword}}},
    description={Fungsi LLM untuk mengekstraksi dan mengklasifikasikan kata kunci},
    sort=LLMkeyword
}

% ============================================
% SIMBOL INFORMASI LOKAL & GLOBAL
% ============================================

\newglossaryentry{sym:Llocal}{
    type=symbols,
    name={\ensuremath{\mathcal{L}_{\text{local}}}},
    description={Informasi lokal hasil ekstraksi dari subgraf},
    sort=Llocal
}

\newglossaryentry{sym:Elocal}{
    type=symbols,
    name={\ensuremath{\mathcal{E}_{\text{local}}}},
    description={Himpunan entitas dari ekstraksi informasi lokal},
    sort=Elocal
}

\newglossaryentry{sym:Rlocal}{
    type=symbols,
    name={\ensuremath{\mathcal{R}_{\text{local}}}},
    description={Himpunan relasi dari ekstraksi informasi lokal},
    sort=Rlocal
}

\newglossaryentry{sym:Tlocal}{
    type=symbols,
    name={\ensuremath{\mathcal{T}_{\text{local}}}},
    description={Himpunan unit teks dari ekstraksi informasi lokal},
    sort=Tlocal
}

\newglossaryentry{sym:pipelinelocal}{
    type=symbols,
    name={\ensuremath{\text{pipeline}_{\text{local}}}},
    description={Fungsi \textit{pipeline} untuk ekstraksi informasi lokal},
    sort=pipelinelocal
}

\newglossaryentry{sym:Lglobal}{
    type=symbols,
    name={\ensuremath{\mathcal{L}_{\text{global}}}},
    description={Informasi global hasil ekstraksi dari jaringan relasi},
    sort=Lglobal
}

\newglossaryentry{sym:Eglobal}{
    type=symbols,
    name={\ensuremath{\mathcal{E}_{\text{global}}}},
    description={Himpunan entitas dari ekstraksi informasi global},
    sort=Eglobal
}

\newglossaryentry{sym:Rglobal}{
    type=symbols,
    name={\ensuremath{\mathcal{R}_{\text{global}}}},
    description={Himpunan relasi dari ekstraksi informasi global},
    sort=Rglobal
}

\newglossaryentry{sym:Tglobal}{
    type=symbols,
    name={\ensuremath{\mathcal{T}_{\text{global}}}},
    description={Himpunan unit teks dari ekstraksi informasi global},
    sort=Tglobal
}

\newglossaryentry{sym:pipelineglobal}{
    type=symbols,
    name={\ensuremath{\text{pipeline}_{\text{global}}}},
    description={Fungsi \textit{pipeline} untuk ekstraksi informasi global},
    sort=pipelineglobal
}

% ============================================
% SIMBOL FUNGSI & OPERATOR
% ============================================

\newglossaryentry{sym:fIE}{
    type=symbols,
    name={\ensuremath{f_{IE}}},
    description={Fungsi pemetaan \textit{Information Extraction} dari teks ke representasi terstruktur},
    sort=fIE
}

\newglossaryentry{sym:argmax}{
    type=symbols,
    name={\ensuremath{\arg\max}},
    description={Operator untuk menemukan argumen yang memaksimalkan suatu fungsi},
    sort=argmax
}

\newglossaryentry{sym:sum}{
    type=symbols,
    name={\ensuremath{\sum}},
    description={Operator penjumlahan},
    sort=sum
}

\newglossaryentry{sym:subset}{
    type=symbols,
    name={\ensuremath{\subseteq}},
    description={Relasi himpunan bagian (\textit{subset})},
    sort=subset
}

\newglossaryentry{sym:in}{
    type=symbols,
    name={\ensuremath{\in}},
    description={Relasi keanggotaan himpunan},
    sort=in
}

\newglossaryentry{sym:mid}{
    type=symbols,
    name={\ensuremath{\mid}},
    description={Simbol kondisional (``diberikan'') dalam probabilitas},
    sort=mid
}

% ============================================
% SIMBOL NUMERIK & THRESHOLD
% ============================================

\newglossaryentry{sym:lessthan50}{
    type=symbols,
    name={\ensuremath{<50}},
    description={Threshold minimum jumlah kata untuk abstrak yang diproses},
    sort=50
}

\newglossaryentry{sym:2024chars}{
    type=symbols,
    name={2.024 karakter},
    description={Ukuran \textit{chunk} untuk \textit{text chunking} dalam \textit{vector indexing}},
    sort=2024
}

\newglossaryentry{sym:50overlap}{
    type=symbols,
    name={50 karakter},
    description={Ukuran \textit{overlap} antar \textit{chunk} untuk menjaga kontinuitas konteks},
    sort=50overlap
}

% ============================================
% CATATAN PENGGUNAAN
% ============================================
% 
% Untuk menggunakan simbol dalam teks, gunakan perintah:
%   \gls{sym:namasimbol}
%
% Contoh:
%   Graf pengetahuan \gls{sym:G} terdiri dari himpunan entitas \gls{sym:E}...
%   Sistem mencari jawaban optimal \gls{sym:astar} berdasarkan \gls{sym:q}...
%
% Untuk menampilkan daftar simbol, gunakan:
%   \printglossary[type=symbols,title={Daftar Simbol}]
%
% Pastikan di preamble dokumen sudah ada:
%   \usepackage[symbols,nogroupskip]{glossaries-extra}
%   \makeglossaries
%   \newglossary[slg]{symbols}{syi}{syg}{Daftar Simbol}
%
